% =============================================================================
% SECTION 1: INTRODUCTION
% =============================================================================
\section{Introduction}\label{sec:intro}

The past three years have witnessed the emergence of \emph{foundation} machine learning interatomic potentials (MLIPs): models pretrained on chemically diverse density functional theory (DFT) corpora at scale, designed to generalize across the periodic table without element-specific retraining\cite{batatia2024mace, deng2023chgnet, yang2024mattersim, barroso2024omat24}. Representative examples include MACE-MP-0 and its successor MACE-MPA-0\cite{batatia2024mace}, CHGNet and its r$^2$SCAN variant\cite{deng2023chgnet}, the Orb family (Orb-v2, Orb-v3)\cite{qian2023orb}, SevenNet\cite{park2024sevennet}, EquiformerV2 (eqV2)\cite{liao2024equiformerv2}, GNoME\cite{merchant2023gnome}, MatterSim\cite{yang2024mattersim}, eSEN\cite{om枯萎}, PET-MAD\cite{mazitov2025petmad}, GRACE\cite{bochkarev2024grace}, DPA3\cite{zhang2024dpa3}, and the OMat24/UMA lineage\cite{barroso2024omat24}. These models share a common ambition: to provide a single potential energy surface (PES) surrogate applicable to arbitrary combinations of elements, phases, and thermodynamic conditions.

Despite this shared ambition, the empirical literature documents striking \emph{heterogeneity} in their error structures. The Matbench Discovery leaderboard\cite{riebesell2025matbench} ranks 47 distinct models on a unified WBM test set with DFT reference convex hull, spanning at least eleven architectural families. The headline F1 scores for stability classification vary by more than a factor of three across compliant models, and the $\kappa_{\text{SRME}}$ thermal-conductivity benchmark\cite{pota2024kappa} produces a different rank ordering than the energy-MAE ordering. MLIP Arena\cite{chiang2025mliparena} extends this critique to physics-aware properties---pairwise potential energy curve (PEC) accuracy, molecular dynamics stability under NVT ramps from 300~K to 3000~K, and chemical reactivity---finding that ``many top-ranked models on bulk crystal equation of state and Matbench Discovery perform poorly in pairwise interactions, suggesting that apparent benchmark success may result from plausible many-body error cancellation in bulk systems.''\cite{chiang2025mliparena} The PET-MAD study\cite{mazitov2025petmad} demonstrates explicitly that ``recent `universal' models deliver qualitative accuracy across the periodic table but are often biased toward low-energy configurations,'' and that this bias diminishes when the training distribution is broadened. Deng \emph{et al.}\cite{deng2025systematic} document ``a consistent PES softening effect in three uMLIPs: M3GNet, CHGNet, and MACE-MP-0,'' characterized by energy and force underprediction across surfaces, defects, solid-solution energetics, ion migration barriers, and phonon vibration modes.

The central question of this manuscript is: \textbf{Is this heterogeneous error structure model-specific noise, or does it possess shared class-level geometry?} We prove that the latter is true. Under three definitional conditions on an architectural class $\cF$ of foundation MLIPs, every model $M \in \cF$ admits a per-model error manifold $\cM(M)$ whose structural properties are \emph{class-uniform}: bounded intrinsic dimension, Vandermonde-type Fisher spectrum decay with a rate $\rho(\cF)$ independent of the specific model, sample complexity $O(d \log N)$ for manifold recovery, and stable evolution under generational succession. The theorem has six clauses; clauses (i)--(iii) and (v) are proved rigorously, while clauses (iv) and (vi) are pre-registered predictions with explicit falsification thresholds.

\subsection{Positioning and Scope}

This manuscript is the second of a planned trilogy. The first paper (CMET, currently in revision) provides the manifold-extension primitive: given a single foundation MLIP $M$, it proves that the set of atomic configurations on which $M$'s prediction error exceeds a threshold forms a smooth submanifold of the configuration space under suitable regularity conditions. The present paper promotes this single-model statement to a \emph{class-level} universality theorem, asserting that the geometric properties of these per-model error manifolds are uniform across all $M \in \cF$. The third paper (forthcoming) will treat the \emph{refusal-mode} behavior of foundation MLIPs---the conditions under which a model should abstain from prediction rather than emit an unreliable energy or force estimate.

The architectural class $\cF$ is defined precisely in Section~\ref{sec:class-f}. Briefly, condition (F1) requires pretraining on a chemically diverse DFT corpus at scale (Materials Project trajectory data, Alexandria/sAlex, OMat24, or MAD); condition (F2) requires representational expressivity strictly above the Pozdnyakov--Ceriotti incompleteness floor\cite{pozdnyakov2020incompleteness}, operationalized via the geometric Weisfeiler--Leman hierarchy of Joshi \emph{et al.}\cite{joshi2023expressive}; and condition (F3) requires smoothness of $M$ as a function of atomic positions outside a measure-zero set of singular configurations, anchored to the smooth basis construction of Bigi \emph{et al.}\cite{bigi2022smooth}. Models satisfying (F1)--(F3) include MACE-MP-0/MPA-0, CHGNet, Orb-v3, MatterSim, eqV2, PET-MAD, SevenNet, GRACE, eSEN, and DPA3; counterexamples include classical empirical potentials (fail F2), the original GAP without equivariant features (fails F2 at scale), and incomplete-descriptor models (fail F2).

\subsection{Summary of Results}

The Universality Theorem (Section~\ref{sec:theorem}) states six clauses:
\begin{enumerate}[label=(\roman*)]
    \item \textbf{Bounded intrinsic dimension.} The per-model error manifold $\cM(M)$ has intrinsic dimension $d(M) \leq d(\cF)$, where $d(\cF)$ depends only on the class constants in (F1)--(F3).
    \item \textbf{Sample complexity.} Recovery of $\cM(M)$ from $n$ residual samples succeeds with probability $1-\delta$ provided $n \geq C \, d(\cF) \log(N/d(\cF)) + C' \log(1/\delta)$, where $N$ is the ambient configuration-space dimension.
    \item \textbf{Cross-model Vandermonde decay.} The empirical Fisher information matrix $F_n(M)$ exhibits geometric eigenvalue decay: $\sigma_m(F_n(M)) \leq \sigma_1(F_n(M)) \cdot \exp(-\rho(\cF)(m-1)) + O_p(\sqrt{\log d / n})$, uniformly over $M \in \cF$. The decay rate $\rho(\cF)$ is a class constant; this is the headline mathematical result.
    \item \textbf{Active learning excess risk.} (Pre-registered prediction.) For active learning with $T$ queries, the excess risk on the residual manifold scales as $\epsilon(T) = C \cdot d(\cF) \log N / T$, with $C$ within a factor of 4 of the Raj--Bach\cite{raj2022convergence} constant.
    \item \textbf{Two-mode inference.} The configuration space partitions into a prediction region (where $M$'s nearest-point projection onto $\cM(M)$ is well-defined and Lipschitz) and a refusal region (where the reach condition fails), with the boundary characterized by Federer's reach machinery\cite{federer1959curvature}.
    \item \textbf{Generational stability.} (Pre-registered prediction.) For consecutive generations $M_g, M_{g+1}$ of a foundation MLIP family, the top-5 principal directions of the cross-configuration residual matrix have pairwise cosine similarity $\geq 0.7$.
\end{enumerate}

The Cross-Model Vandermonde Lemma (Section~\ref{sec:vandermonde-proof}) assembles four published primitives---Waterfall \emph{et al.}'s single-model Vandermonde bound\cite{waterfall2006sloppy}, Tropp's matrix Bernstein inequality\cite{tropp2012user}, Wedin's $\sin\theta$ theorem\cite{wedin1972perturbation} / Dopico's refinement\cite{dopico2000sin}, and Beckermann's Vandermonde conditioning lower bound\cite{beckermann2000condition}---into a class-uniform statement. The $\Delta$-learning corollary (Section~\ref{sec:escape}) gives an upper bound on residual error after correction by a Gaussian process regression on the base MLIP's embedding, in terms of the projection of the residual onto $\cM(M)$ and the dimension $d(\cF)$.

Two predictions are pre-registered with explicit falsification thresholds (Section~\ref{sec:predictions}). Prediction P1 ties the active-learning excess risk to the Raj--Bach non-asymptotic rate; falsification requires either the constant $C$ or the scaling exponent in $T$ to deviate by more than the stated factor. Prediction P2 ties generational stability to the Matbench Discovery WBM test set; falsification requires cosine similarity below 0.5 for any consecutive generational transition within $\cF$.
